\documentclass[10pt,a4paper]{amsart}
\usepackage[margin=19mm]{geometry}
\usepackage{amsmath,amssymb,mathtools}
\usepackage[scr=rsfs,cal=euler]{mathalfa}
\usepackage{xfrac}
\usepackage[unicode,hidelinks,bookmarks=false]{hyperref}
\newcommand{\ie}{i.e.\ }
\newcommand{\cf}{cf.\ }
\newcommand{\resp}{respectively\ }
\newcommand{\Subsection}{\subsection}
\providecommand{\nobreakdash}{\nobreak\hbox{-}}
\setlength{\emergencystretch}{2em}
\multlinegap0pt
\usepackage{csquotes}
\usepackage{xcolor}
\usepackage{amssymb}
\usepackage{stackengine}
\usepackage[shortlabels]{enumitem}
\usepackage{float}
\usepackage{tikz}
\usepackage{tcolorbox}
\newtheorem{lemma}{Lemma}
\title{\uppercase{Considering The\\ Satisfiability of Cubic\\ Diophantine Equations}}
\author{Milan Rosko}
\date{Source: arXiv:2510.00759v8 (CC BY 4.0)}
\begin{document}
\maketitle

\noindent\emph{Source and licence.} \uppercase{Considering The\\ Satisfiability of Cubic\\ Diophantine Equations}. Version arXiv:2510.00759v8 (CC BY 4.0), distributed by arXiv under the Creative Commons Attribution 4.0 International licence, http://creativecommons.org/licenses/by/4.0/, as recorded in the arXiv licence metadata for this exact version and shown on the abstract page https://arxiv.org/abs/2510.00759v8. Full licence notice: \texttt{licenses/arxiv-2510.00759-CC-BY-4.0.txt}.\par\smallskip

\noindent\textit{Editorial scope.} Counted unit: the article's lemma lem:quartic-cubic-system-closure (source0.tex lines 788--795). The article's complete original proof of it is delivered with it (source0.tex lines 797--818, one \texttt{proof} environment of 22 lines). No mathematics is written, repaired or re-derived: the statement and the proof are the article's own lines, in the article's own notation, and no retained line is substituted or re-flowed.  No further source-local context is needed: every cross-reference of every retained line resolves inside the two delivered spans.  Nothing was omitted from any retained span, so the package carries no omission record.  No retained line contains a citation, so the wrapper carries no bibliography and no bibliography entry.  No retained line contains a figure environment, an image-inclusion command or drawing code, so no image is delivered and the package declares no asset dependency.  Editorial formatting only, in addition: an AMS \texttt{amsart} preamble that restates the article's own declarations and macros used by the retained lines, namely the environments \texttt{lemma} on separate counters, together with the standard packages those lines need.  The retained environments print in this document as Lemma 1 in order of appearance, whereas the article prints the same items with the numbering of its own full document.  The complete native source of the article as distributed in the arXiv e-print for this exact version is bundled unchanged as \texttt{sources/186157/source0.tex}.
\begin{lemma}[Quartic-to-cubic systems]
\label{lem:quartic-cubic-system-closure}
	Every quartic Diophantine equation in $v$ variables can be transformed, uniformly and effectively, into a finite system of Diophantine equations of degree at most $3$ in
	\begin{equation}
		v+\binom{v+1}{2}
	\end{equation}
	variables.
\end{lemma}

\begin{proof}[Proof sketch]
	Let the original variables be $x_1,\ldots,x_v$. For every pair
	$1\le i\le j\le v$, introduce a new variable $p_{ij}$ intended to represent the product $x_i x_j$, and add the quadratic equation
	\begin{equation}
		p_{ij}=x_i x_j.
	\end{equation}
	Every quartic monomial
	\begin{equation}
		x_a x_b x_c x_d
	\end{equation}
	can then be rewritten as a product of two quadratic product variables, for example
	\begin{equation}
		x_a x_b x_c x_d
		\quad\longmapsto\quad
		p_{ab}p_{cd},
	\end{equation}
	after fixing a uniform convention for ordering the pairs.

	The rewritten main equation has degree at most $2$ in the enlarged variables, while the defining equations $p_{ij}=x_i x_j$ have degree $2$. Cubic and lower degree terms of the original equation are left unchanged. Hence the output is a finite system of equations of degree at most $3$.

	Solvability is preserved. Any solution of the original quartic equation extends to a solution of the system by setting $p_{ij}=x_i x_j$. Conversely, any solution of the system satisfies the original quartic equation after substituting back the enforced products.
\end{proof}

\end{document}
